\documentclass[11pt]{scrartcl}

\usepackage[stixtwo]{fontsetup}
\usepackage{amsmath}
\usepackage{tikz}

\begin{document}
\begin{enumerate}
	\item \[
		      \frac{5 + 2i}{i + 3}
		      = \frac{(5 + 2i)(3 - i)}{(3 + i)(3 - i)}
		      = \frac{15 - 5i + 6i - 2i^2}{10}
		      = \frac{17}{10} + \frac{1}{10} i
	      \]

	\item $-27i = 27\left(\cos \frac{3\pi}{2} + i \sin \frac{3\pi}{2}\right)$,
	      így
	      \[
		      z_k
		      = 3\left(
		      \cos\left(\frac{\pi}{2} + \frac{2k\pi}{3}\right)
		      + i \sin\left(\frac{\pi}{2} + \frac{2k\pi}{3}\right)
		      \right),
		      \quad k = 0, 1, 2,
	      \]
	      azaz $z_0 = 3e^{i \frac{\pi}{2}} = 3i$, $z_1 = 3e^{i \frac{7\pi}{6}}$
	      és $z_2 = 3e^{i \frac{11\pi}{6}}$. A gyökök egy origó középpontú,
	      $3$ sugarú körbe írt szabályos háromszög csúcsai:
	      \begin{center}
		      \begin{tikzpicture}[scale=0.8, >=stealth]
			      \draw[->] (-2.8, 0) -- (2.8, 0) node[right] {$\operatorname{Re}$};
			      \draw[->] (0, -2.8) -- (0, 2.8) node[above] {$\operatorname{Im}$};
			      \draw[dashed] (0, 0) circle (2);
			      \node[below right] at (2, 0) {$3$};
			      \draw[thin] (90:2) -- (210:2) -- (330:2) -- cycle;
			      \fill (90:2) circle (2pt) node at (102:2.5) {$z_0$};
			      \fill (210:2) circle (2pt) node at (222:2.5) {$z_1$};
			      \fill (330:2) circle (2pt) node at (342:2.5) {$z_2$};
		      \end{tikzpicture}
	      \end{center}
\end{enumerate}
\end{document}
